\section*{Contoh}
\addcontentsline{toc}{section}{Contoh}

\begin{contoh}[Pembuktian dengan induksi]
Buktikan bahwa untuk setiap \(n \ge 1\),
\[
\sum_{i=1}^{n} i = \frac{n(n+1)}{2}.
\]
\textbf{Basis.} Untuk \(n=1\), ruas kiri bernilai \(1\) dan ruas
kanan bernilai \(1(2)/2 = 1\). Basis benar.

\textbf{Langkah induksi.} Asumsikan sifat benar untuk \(k\):
\[
\sum_{i=1}^{k} i = \frac{k(k+1)}{2}.
\]
Maka
\[
\sum_{i=1}^{k+1} i
= \frac{k(k+1)}{2} + (k+1)
= \frac{k(k+1) + 2(k+1)}{2}
= \frac{(k+1)(k+2)}{2}.
\]
Bentuk terakhir adalah rumus untuk \(n = k+1\). Dengan induksi,
sifat terbukti untuk semua \(n \ge 1\).
\end{contoh}

\begin{contoh}[Membuktikan bahasa]
Misalkan \(G = (\{S\}, \{a,b\}, P, S)\) dengan aturan
\(S \to aSb \mid \varepsilon\). Tunjukkan bahwa
\(L(G) = \{a^n b^n \mid n \ge 0\}\).

\textbf{Arah pertama.} Setiap derivasi yang mungkin adalah
\[
S \Rightarrow aSb \Rightarrow a^2 S b^2 \Rightarrow \cdots
\Rightarrow a^n S b^n \Rightarrow a^n b^n.
\]
Dengan induksi pada jumlah langkah, setiap string yang diturunkan
berbentuk \(a^n b^n\).

\textbf{Arah kedua.} Untuk setiap \(n \ge 0\), string \(a^n b^n\)
dapat diturunkan dengan menerapkan aturan \(S \to aSb\) sebanyak
\(n\) kali, lalu \(S \to \varepsilon\). Jadi setiap string berbentuk
tersebut berada di \(L(G)\).

Kedua arah membuktikan \(L(G) = \{a^n b^n \mid n \ge 0\}\).
\end{contoh}

\begin{contoh}[Derivasi dan pohon urai]
Diberikan tata bahasa
\[
E \to E + T \mid T, \quad
T \to T * F \mid F, \quad
F \to (E) \mid \code{id}.
\]
Tentukan derivasi leftmost dan pohon urai untuk
\(\code{id} * (\code{id} + \code{id})\).

Derivasi leftmost:
\[
\begin{aligned}
E &\Rightarrow T\\
  &\Rightarrow T * F\\
  &\Rightarrow F * F\\
  &\Rightarrow \code{id} * F\\
  &\Rightarrow \code{id} * (E)\\
  &\Rightarrow \code{id} * (E + T)\\
  &\Rightarrow \code{id} * (T + T)\\
  &\Rightarrow \code{id} * (F + T)\\
  &\Rightarrow \code{id} * (\code{id} + T)\\
  &\Rightarrow \code{id} * (\code{id} + F)\\
  &\Rightarrow \code{id} * (\code{id} + \code{id}).
\end{aligned}
\]
Pohon uraiya memiliki akar \(E\), cabang kiri \(T\) yang menurun ke
\(\code{id}\), operator \(*\), dan cabang kanan \(F\) yang menurun ke
\((E)\). Di dalam kurung, ekspresi \(E + T\) diurai dengan cara yang
sama seperti contoh sebelumnya.
\end{contoh}

\begin{contoh}[Klasifikasi bahasa]
Tentukan kelas bahasa berikut pada hierarki Chomsky.
\begin{enumerate}[label=(\alph*),leftmargin=*]
  \item \(L_1 = \{w \in \{a,b\}^{*} \mid w \text{ memuat } aba\}\).
  \item \(L_2 = \{a^n b^n \mid n \ge 0\}\).
  \item \(L_3 = \{a^n b^n c^n \mid n \ge 0\}\).
\end{enumerate}
\textbf{Jawaban.}
\begin{enumerate}[label=(\alph*),leftmargin=*]
  \item \(L_1\) regular. Bahasa ini dikenali oleh DFA dengan empat
        state yang mengingat apakah prefiks terakhir adalah
        \(\varepsilon\), \(a\), \(ab\), atau \(aba\).
  \item \(L_2\) bebas konteks, tetapi bukan regular. Tata bahasa
        \(S \to aSb \mid \varepsilon\) membangkitkannya. Bahasa ini
        memerlukan penghitungan jumlah \(a\) dan \(b\) yang tidak
        terbatas, sehingga DFA dengan state berhingga tidak cukup.
  \item \(L_3\) sensitif konteks, tetapi bukan bebas konteks. Bahasa
        ini memerlukan pencocokan tiga kelompok simbol, yang tidak
        dapat dilakukan oleh satu stack.
\end{enumerate}
\end{contoh}
